% Conclusions and Recommendations (BGU \S 10.7.2.9).
% V6: plain-language reading woven into the four conclusions; Goal 3 and the
% calibration recommendation corrected against the frozen data (TransNeXt is
% ~6.5x ResNet50, not 1.5x; ECE does not grow toward L5); the plan-versus-
% execution ledger rebuilt from the PRD v5 section 4.9 numbers (the V5 table
% understated the B2/C2 actuals).
\section{Conclusions and Recommendations}
\label{sec:conclusions}

\subsection{Scientific conclusions}
This work delivered a reproducible 402-cell training campaign across three
architectures (ResNet50, DenseNet121, TransNeXt-tiny), two datasets
(CIFAR-10, MNIST), five degradation levels, and seven experimental phases,
anchored to a deterministic per-pipeline PSNR / SSIM probe and a 48-replicate
multi-seed audit. Four conclusions follow from the measurements.

\begin{enumerate}
  \item \textbf{The collapse is an information bottleneck, not
        over-fitting.} The strongest regularization treatment recovers less
        than one percentage point of the combined-degradation collapse
        (Phase~D: +0.77~\Mpp{} on average), and moving to the simpler THz
        protocol recovers about four (Phase~B2: +4.01~\Mpp{}, roughly three
        quarters of which is the protocol itself rather than the
        regularization). If the models were merely over-fitting the degraded
        distribution, training-side intervention would have closed a large
        part of the gap; it does not. The accuracy is gone because the
        information is gone.
  \item \textbf{The foveal transformer is the most robust backbone at the
        operationally relevant levels.} TransNeXt-tiny leads the CNNs by
        about ten \Mpp{} at the mild-to-moderate end of the combined curve on
        CIFAR-10 and is the most robust model on every single-axis
        comparison at matched image quality. The advantage is real (well
        above the seed noise) but conditional: it exists exactly where the
        degraded image still contains texture for its attention to use.
  \item \textbf{At extreme degradation, architecture stops mattering.} At
        Phase~B1 L5 all three backbones land in a 19--21\% band on CIFAR-10
        (27--28\% on MNIST) --- within seed noise of one another. Once the
        spatial information is destroyed, no backbone in this design space,
        and plausibly none outside it, can recover what is no longer there.
  \item \textbf{Resolution is the dominant degradation axis.} On the
        measured PSNR and SSIM scales, the resolution axis costs roughly
        four times more accuracy than blur or salt-and-pepper
        ($-27.45$~\Mpp{} mean L1--L5 against $-6.58$ and $-5.13$) and shows
        an almost perfect accuracy--quality association
        ($\rho = +0.98$). For a THz system engineer this is the actionable
        result: spatial resolution is where the accuracy goes, so spatial
        resolution is what front-end investment should restore first.
\end{enumerate}

\subsection{Project evaluation against the preliminary goals}
The preliminary design report (PDR) set three acceptance criteria.

\begin{itemize}
  \item \textbf{Goal 1 --- at least 80\% top-1 accuracy at the moderate (L3)
        operating point.} \emph{Partially met, and informatively so.} On the
        full combined L3 pipeline the target is missed on CIFAR-10
        (36--39\%), essentially met on MNIST (79--81\%), and comfortably
        exceeded on the milder single-axis operating points (for example,
        every Phase~C2 blur and salt-and-pepper cell at L3). The CIFAR-10
        shortfall is not an engineering failure but the central scientific
        finding: the moderate combined degradation destroys enough
        information that \emph{no} backbone in the design space reaches
        80\% --- and Phase~D shows that no training trick changes this.
  \item \textbf{Goal 2 --- an identifiable architectural winner beating the
        baseline by a meaningful margin.} \emph{Met, in the regime where a
        winner can exist.} TransNeXt-tiny beats the CNN mean by about ten
        \Mpp{} at L1--L2 on CIFAR-10, far above the multi-seed noise floor
        ($\sigma < 1.2$~\Mpp{}). The margin decays to about two points at L3
        and vanishes by L4, so the ``winner'' claim is explicitly scoped to
        the mild-to-moderate range where it is decision-relevant.
  \item \textbf{Goal 3 --- inference latency compatible with sensor-side
        deployment.} \emph{Met for all three backbones, with very different
        margins.} At batch size 1 on the RTX 5070: ResNet50 4.2~ms
        (238 images/s), DenseNet121 10.7~ms (94 images/s), TransNeXt-tiny
        27.5~ms (36 images/s). All three clear the frame rate of practical
        THz scanners; but the transformer is roughly $6.5\times$ slower than
        ResNet50 because its foveal-attention path is not graph-compiled on
        the current stack, leaving it the thinnest deployment margin ---
        the subject of a continuation recommendation below.
\end{itemize}

\textbf{Trade-offs of the final architecture choice.}\quad ResNet50 is the
fastest and simplest backbone and a solid baseline, but it carries no
robustness advantage. DenseNet121 adds a small edge at moderate CIFAR-10
degradation through feature reuse, at the cost of $2.5\times$ ResNet50's
latency and higher memory. TransNeXt-tiny delivers the best accuracy at mild
and moderate degradation, at the cost of $6.5\times$ ResNet50's latency and
no advantage at the extreme end. For a THz deployment operating at
mild-to-moderate degradation the recommendation is TransNeXt-tiny; at
extreme degradation the engineering budget is better spent on front-end
restoration than on any classifier.

\subsection{Plan versus execution}
Table~\ref{tab:plan_exec} compares the planned compute budget with the
logged execution from the campaign ledger. The overall envelope closed at
about +9\% --- comfortably inside the $\sim$310 GPU-hour saving created by
the \texttt{transnext\_small} $\rightarrow$ \texttt{tiny} substitution ---
but the per-phase picture is uneven and worth stating honestly: the
convergence-first policy (60-epoch budget, early-stopping patience 10) let
the T3-regularized Phase~B2 / C2 cells train materially longer than
estimated (+45\% and +48\% over their line items), while the 186-cell
baseline re-run and the post-training diagnostics came in well under. All
scheduled milestones --- PDR, the determinism gates, the 402-cell campaign,
the multi-seed audit, the diagnostics, and this report --- were met against
the 26 July 2026 submission deadline.

\begin{table}[htbp]
\centering
\caption{Plan versus execution (GPU-hours). Planned: PDR / PRD estimates; actual: campaign ledger; $\approx$ marks ledger estimates.}
\label{tab:plan_exec}
\setlength{\tabcolsep}{5pt}
\renewcommand{\arraystretch}{1.15}
\begin{tabular}{lrrl}
\toprule
Work item & Planned & Actual & Status \\
\midrule
Hyperparameter pre-tuning (Optuna)        & 20.0  & $\approx$20 & Within budget \\
Phases A / B1 / C (186 cells)             & 62.0  & $\approx$50 & $\sim$19\% under \\
Phase D regularization sweep (90 cells)   & 52.5  & 56.5        & $+7.6\%$ over \\
Phase B2 sweep (30 cells)                 & 15.0  & 21.7        & $+45\%$ over \\
Phase B2-nr arm (6 cells)                 & 3.5   & 3.8         & Within budget \\
Phase C2 sweep (90 cells)                 & 48.0  & 70.9        & $+48\%$ over \\
Multi-seed audit (48 replicates)          & 24.0  & $\approx$25 & Within budget \\
Post-training diagnostics (no re-train)   & 3.5   & 0.7         & $\sim$80\% under \\
\midrule
Total                                     & 228.5 & $\approx$249 & $+9\%$ \\
\midrule
\multicolumn{4}{l}{\emph{Architecture substitution (\texttt{small}$\rightarrow$\texttt{tiny}): $\approx -310$ GPU-h vs.\ the original plan.}} \\
\bottomrule
\end{tabular}
\end{table}

\subsection{Limitations}
The conclusions are bounded by the experimental scope. The five-axis pipeline
is a synthetic mathematical model: it omits wavelength-dependent diffraction,
mixing-element nonlinearity, and the dead-pixel patterns of real THz
detectors. CIFAR-10 and MNIST are small benchmarks, so the per-axis
attribution is not validated at ImageNet scale or on true THz object
categories. The Phase~B2 / C2 sweeps were run under the T3 treatment only, so
protocol and regularization interact in all but the L3 cells (where the
B2-nr arm separates them). No learned super-resolution baseline was run (a
scoped-out Phase~E), so the ``restore resolution first'' recommendation has
not yet been tested against an actual restoration pipeline. Finally, the
acceptance criteria are model-side only; no human-perception bound was
collected for the extreme operating points.

\subsection{Recommendations for continuation}
\begin{enumerate}
  \item \textbf{Acquire real THz data.}\quad Even a small dataset at the
        target frequency would calibrate the synthetic model against ground
        truth and re-weight the per-axis attribution accordingly.
  \item \textbf{Add a super-resolution pre-stage.}\quad Resolution dominates
        the attribution by $4\times$, so the highest-value next ablation
        places a learned upsampler in front of each backbone and re-runs
        Phases~B2 / C2. If the attribution is right, this single stage
        should recover more accuracy than any classifier change measured
        here.
  \item \textbf{Graph-compile the TransNeXt attention path.}\quad The
        foveal-attention custom op is the reason for the $6.5\times$ latency
        gap to ResNet50. Compiling it (or replacing it with a fused
        implementation) would widen the transformer's deployment margin at
        no accuracy cost.
  \item \textbf{Reproduce the attribution at scale.}\quad Repeat the
        per-axis analysis on a large robustness benchmark such as
        ImageNet-A [6] to test whether resolution dominance
        holds beyond small benchmarks.
  \item \textbf{Calibration monitoring.}\quad Miscalibration stays moderate
        (mean ECE 5.6\% at L3, 2.9\% at L5) but is not free: the worst cell
        (12.7\%) sits on the unregularized B2-nr arm. A cheap post-hoc
        temperature-scaling step, fitted per operating point, would tighten
        the confidence outputs without touching top-1 accuracy --- worth
        adding before any operational use of the confidence scores.
\end{enumerate}
